% GENERATED FILE. Edit canonical lessons, then run npm run generate:books.
% canonical-source-hash: fnv1a64:9479f29fa42de7a5
% canonical-lessons: RU-C59-rain, RU-C59-snow, RU-C59-wind, RU-C59-weather, RU-C59-world

\chapter{Rain, Snow, Wind, Weather, World}
\label{ch:ru-rain-snow-wind-weather-world}
\hlchaptermodality{59}

\begin{chapteropening}
\textbf{By the end of this chapter:} \emph{I can name rain, snow, the wind and the weather, and the world.}

Complete the last lesson of chapter 59: I can name rain, snow, the wind and the weather, and the world.
\end{chapteropening}

\section[\texorpdfstring{\ru{дождь} (dozhd)}{дождь (dozhd)}]{\ru{дождь} (dozhd) — rain}

\label{lesson:RU-C59-rain}



\begin{quote}
Before the new one: say the Russian for a star, then the Russian for a cloud.
\end{quote}

\subsection*{You'll want to know: \ru{дождь}}
\textbf{\ru{дождь}} (\emph{dozhd}) — \textquotedblleft{}rain\textquotedblright{}.

\begin{cousinweb}
\textbf{\ru{дождь}} $\leftarrow$ Proto-Slavic \emph{*dŭždžĭ} — inherited. Vasmer thinks it probably goes with Norwegian \emph{dusk-regn} and Swedish \emph{regn-dusk}, \textquotedblleft{}drizzle\textquotedblright{}; he also reports a view that starts from \emph{*dus-dius}, \textquotedblleft{}bad weather\textquotedblright{}, so the origin is debated.
\end{cousinweb}

\begin{grammarlens}[title={What you've built}]
One more word for this chapter.
\end{grammarlens}

\subsection*{Your turn}
\begin{itemize}
\raggedright
  \item \emph{Say it:} \emph{dozhd}
  \item \emph{Say it:} \emph{dozhd}, once more
  \item \emph{Say it:} say \emph{dozhd}
  \item \emph{Recall:} say the Russian for a star, then the Russian for a cloud, then say \emph{dozhd} again
\end{itemize}

\begin{tcolorbox}[breakable,colback=teal!4,colframe=teal!35!black,title={Before you move on}]
What does \ru{дождь} mean? (\textquotedblleft{}Rain\textquotedblright{}.) Say it once more.
\end{tcolorbox}
\section[\texorpdfstring{\ru{снег} (sneg)}{снег (sneg)}]{\ru{снег} (sneg) — snow}

\label{lesson:RU-C59-snow}



\begin{quote}
Before the new one: say the Russian for a cloud, then the Russian for rain.
\end{quote}

\subsection*{You'll want to know: \ru{снег}}
\textbf{\ru{снег}} (\emph{sneg}) — \textquotedblleft{}snow\textquotedblright{}.

\begin{cousinweb}
\textbf{\ru{снег}} $\leftarrow$ Proto-Slavic \emph{*sněgŭ} — inherited. Vasmer lists Lithuanian \emph{sniẽgas}, Gothic \emph{snaiws}, Greek \emph{nípha} and Latin \emph{nix}, all \textquotedblleft{}snow\textquotedblright{}: English \emph{snow} is a cousin, from PIE \emph{*sneyg\textsuperscript{w}\textsuperscript{h}-}, \textquotedblleft{}to snow\textquotedblright{}.
\end{cousinweb}

\begin{grammarlens}[title={What you've built}]
One more word for this chapter.
\end{grammarlens}

\subsection*{Your turn}
\begin{itemize}
\raggedright
  \item \emph{Say it:} \emph{sneg}
  \item \emph{Say it:} \emph{sneg}, once more
  \item \emph{Say it:} say \emph{sneg}
  \item \emph{Recall:} say the Russian for a cloud, then the Russian for rain, then say \emph{sneg} again
\end{itemize}

\begin{tcolorbox}[breakable,colback=teal!4,colframe=teal!35!black,title={Before you move on}]
What does \ru{снег} mean? (\textquotedblleft{}Snow\textquotedblright{}.) Say it once more.
\end{tcolorbox}
\section[\texorpdfstring{\ru{ветер} (véter)}{ветер (véter)}]{\ru{ветер} (véter) — the wind}

\label{lesson:RU-C59-wind}



\begin{quote}
Before the new one: say the Russian for rain, then the Russian for snow.
\end{quote}

\subsection*{You'll want to know: \ru{ветер}}
\textbf{\ru{ветер}} (\emph{véter}) — \textquotedblleft{}the wind\textquotedblright{}.

\begin{cousinweb}
\textbf{\ru{ветер}} $\leftarrow$ Proto-Slavic \emph{*větrŭ} — inherited, from the root of \emph{\ru{веять}}, \textquotedblleft{}to blow\textquotedblright{} (Vasmer). Vasmer lists Lithuanian \emph{vėjas}, \textquotedblleft{}wind\textquotedblright{}, and Sanskrit \emph{vāta-}, \textquotedblleft{}wind\textquotedblright{}; English \emph{wind} comes from the same PIE root \emph{*h\textsubscript{2}weh\textsubscript{1}-}, \textquotedblleft{}to blow\textquotedblright{}.
\end{cousinweb}

\begin{grammarlens}[title={What you've built}]
One more word for this chapter.
\end{grammarlens}

\subsection*{Your turn}
\begin{itemize}
\raggedright
  \item \emph{Say it:} \emph{véter}
  \item \emph{Say it:} \emph{véter}, once more
  \item \emph{Say it:} say \emph{véter}
  \item \emph{Recall:} say the Russian for rain, then the Russian for snow, then say \emph{véter} again
\end{itemize}

\begin{tcolorbox}[breakable,colback=teal!4,colframe=teal!35!black,title={Before you move on}]
What does \ru{ветер} mean? (\textquotedblleft{}The wind\textquotedblright{}.) Say it once more.
\end{tcolorbox}
\section[\texorpdfstring{\ru{погода} (pogóda)}{погода (pogóda)}]{\ru{погода} (pogóda) — the weather}

\label{lesson:RU-C59-weather}



\begin{quote}
Before the new one: say the Russian for snow, then the Russian for the wind.
\end{quote}

\subsection*{You'll want to know: \ru{погода}}
\textbf{\ru{погода}} (\emph{pogóda}) — \textquotedblleft{}the weather\textquotedblright{}.

\begin{cousinweb}
\textbf{\ru{погода}} is \emph{\ru{по}-} with \emph{\ru{год}}, which in Old Russian meant \textquotedblleft{}time, a fitting moment\textquotedblright{} (Vasmer); the same root is in \emph{\ru{годиться}}, \textquotedblleft{}to be fit\textquotedblright{}. Weather was first \textquotedblleft{}fair, fitting time\textquotedblright{}, as in \emph{\ru{непогода}}, \textquotedblleft{}bad weather\textquotedblright{}.
\end{cousinweb}

\begin{grammarlens}[title={What you've built}]
One more word for this chapter.
\end{grammarlens}

\subsection*{Your turn}
\begin{itemize}
\raggedright
  \item \emph{Say it:} \emph{pogóda}
  \item \emph{Say it:} \emph{pogóda}, once more
  \item \emph{Say it:} say \emph{pogóda}
  \item \emph{Recall:} say the Russian for snow, then the Russian for the wind, then say \emph{pogóda} again
\end{itemize}

\begin{tcolorbox}[breakable,colback=teal!4,colframe=teal!35!black,title={Before you move on}]
What does \ru{погода} mean? (\textquotedblleft{}The weather\textquotedblright{}.) Say it once more.
\end{tcolorbox}
\section[\texorpdfstring{\ru{мир} (mir)}{мир (mir)}]{\ru{мир} (mir) — the world; peace}

\label{lesson:RU-C59-world}



\begin{quote}
Before the new one: say the Russian for the wind, then the Russian for the weather.
\end{quote}

\subsection*{You'll want to know: \ru{мир}}
\textbf{\ru{мир}} (\emph{mir}) — \textquotedblleft{}the world; peace\textquotedblright{}. Both meanings in one word.

\begin{cousinweb}
\textbf{\ru{мир}} $\leftarrow$ Proto-Slavic \emph{*mirŭ}, \textquotedblleft{}peace\textquotedblright{} — inherited, related to Old Lithuanian \emph{mieras}, \textquotedblleft{}peace\textquotedblright{}, and Latvian \emph{miêrs} (Vasmer). The sense \textquotedblleft{}village community\textquotedblright{} grew from \textquotedblleft{}peaceful community\textquotedblright{}, and from it \textquotedblleft{}the world\textquotedblright{}.
\end{cousinweb}

\begin{grammarlens}[title={What you've built}]
One more word for this chapter.
\end{grammarlens}

\subsection*{Your turn}
\begin{itemize}
\raggedright
  \item \emph{Say it:} \emph{mir}
  \item \emph{Say it:} \emph{mir}, once more
  \item \emph{Say it:} say \emph{mir}
  \item \emph{Recall:} say the Russian for the wind, then the Russian for the weather, then say \emph{mir} again
\end{itemize}

\begin{tcolorbox}[breakable,colback=teal!4,colframe=teal!35!black,title={Before you move on}]
What does \ru{мир} mean? (\textquotedblleft{}The world; peace\textquotedblright{}.) Say it once more.
\end{tcolorbox}
